\author{OTS}
\documentclass[14pt]{../../inc/mybib}

\setincpath{../../inc/}

\usepackage{bible_style}
\usepackage{textmakro}
\graphicspath{{../../assets/images/}}
\usepackage{header}

\newcommand{\Name}{Thomas}

% ensure scrlayer-scrpage has sufficient footheight
\setlength{\footheight}{20.4pt}

\begin{document}
\begin{bibelbox}{SCHL}{Kol}{3:15-17}
    Lasst das Wort des Christus reichlich in euch wohnen in aller Weisheit; lehrt und ermahnt einander und singt mit Psalmen und Lobgesängen und geistlichen Liedern dem Herrn lieblich in eurem Herzen. Und was immer ihr tut in Wort und Werk, das tut alles im Namen des Herrn Jesus und dankt Gott, dem Vater, durch ihn.
\end{bibelbox}
\section{Begrüssung}
Ich möchte auch euch alle recht herzlich zum Gottesdienst begrüssen. Schön, dass ihr gekommen seid, um unserem \herr N zu danken, ihn zu loben und zu preisen.

\betonung{Ingrid und Stefan Beitze aus lassen uns Guatemala recht herzlich grüssen.}

Auch dich \Name{} will ich herzlich begrüssen. Schön, dass du wieder bei uns bist und mit uns Gottesdienst feierst.
\begin{itemize}
    \item[] \beten{}
    \item[] Und anschliessend singen wir zusammen das Lied
    \item[] \lied{Nur durch Christus in mir}
\end{itemize}

\section{Informationen}
\begin{itemize}
    \item \bt[infos]{Bibel und Gebetsabend:} Do, 05.03.2026 20:00 Uhr Bibel und Gebetsabend ; zu Markus 5,35-43 mit Norbert Lieth.
    \item \bt[infos]{Nächster Gottesdienst:} So, 08.03.2026 14:45 Uhr mit Samuel Rindlisbacher.
    \item \bt[infos]{Kollekte:} Die Kollekte wird hinten in der grünen Box gesammelt und wird vollumfänglich für den Bau dieser Gemeinde eingesetzt.
\end{itemize}

\section{ Input }
\begin{spacing}{1.5}
    \begin{block}[Die Entrückung]
        Ein viel diskutiertes Thema über die Endzeit. Entrückung vor der Trübsal? Entrückung in der Mitte der Trübsal? Oder am Ende der Trübsal? Oder dann vielleicht doch keine Entrückung? Es gibt viele Theorien darüber.

        \rot{FOLIE:}\\
        Das ist keine Wolke auf der wir Harfe spielen werden, sondern soll Jesus wiederkommen symbolisieren, so wie er in den Himmel aufgenommen wurde.
        \rot{ENDE FOLIE}\\
        Wir, der Mitternachtsruf und auch ich, vertreten die Theorie der Vorentrückung oder auch Prätribulationismus genannt. Das heisst, wir glauben, dass die Wiedergeborenen Christen vor der 7-jährigen Trübsalszeit entrückt werden. Laut der Bibel ist es sicher, dass eine Entrückung stattfinden wird. In der Bibel werden 8 Entrückungen erwähnt.
        \begin{enumerate}
            \item Henoch (1Mo 5,24)
            \item Elia (2Kö 2,11-12)
            \item Jesus (Apg 1,9-11)
            \item Philippus (Apg 8,39-40)
            \item Paulus (2Kor 12,2-4)
            \item Johannes (Offb 4,1-2)
            \item Die Entrückung der Gemeinde (1Thess 4,16-17)
            \item Zwei Zeugen (Offb 11,11-12)
        \end{enumerate}
        Es ist also nicht etwas Exotisches, sondern etwas, was in der Bibel mehrfach erwähnt wird. Das Wort Entrückung (Englisch rapture, Deutsch Entrücken), kommt im griechischen als harpazo 14-mal im Neuen Testament vor. Es bedeutet: wegnehmen, wegziehen, wegtragen, fortreißen, entrücken.
        Typische Bibelstellen, die dieses Wort enthalten, sind:
        \begin{itemize}
            \item Matthäus 11,12
                  \begin{bibelbox}{SCHL}{Matt}{11:12}
                      Aber von den Tagen Johannes des Täufers an bis jetzt leidet das Reich der Himmel Gewalt, und die Gewalttätigen \betonung{reißen} es an sich.
                  \end{bibelbox}
            \item Apostelgeschichte 8:39
                  \begin{bibelbox}{SCHL}{Apg}{8:39}
                      Als sie aber aus dem Wasser stiegen, \betonung{entrückte} der Geist des Herrn den Philippus, und der Kämmerer sah ihn nicht mehr, sondern ging fröhlich seines Weges.
                  \end{bibelbox}
        \end{itemize}
        Nur mal zwei Beispiele von den 14 Stellen. Die restlichen Bibelstellen seht hier auf der Folie. Also es gibt das Wort Entrückung in der Bibel und es werden acht Entrückungen erwähnt. So liegt es nahe das auch eine Entrückung stattfinden wird.

        Jetzt kommt die Frage auf, wann wird denn diese Entrückung für uns wiedergeborenen Christen stattfinden? Können wir das irgendwie aus der Bibel herauslesen? Oder gibt es etwa einen Zeitpunkt? Diese Frage ist einfach zu beantworten. Nein, es gibt keinen konkreten Zeitpunkt. Wir wissen, dass es geschieht und wie es geschieht. Das lesen wir in 1Thessalonicher 4,16-17
        \begin{bibelbox}{SCHL}{IThes}{4:15-17}
            Denn das sagen wir euch in einem Wort des Herrn: Wir, die wir leben und bis zum Kommen des Herrn übrig bleiben, werden den Entschlafenen nicht zuvorkommen; denn der Herr selbst wird, wenn der Befehl ergeht und die Stimme des Erzengels und die Posaune Gottes erschallt, vom Himmel herniederfahren, und die Toten in Christus werden zuerst auferstehen. Danach werden wir, die wir leben und übrigbleiben, zugleich mit ihnen entrückt werden in den Wolken, zur Begegnung mit dem Herrn, in die Luft, und so werden wir bei dem Herrn sein allezeit.
        \end{bibelbox}
        So wird es geschehen. Der Herr wird uns entgegenkommen und wir werden mit ihm auf oder in den Wolken entrückt und das wird im Nu geschehen, also plötzlich. Wir können uns nicht darauf verlassen, dass zuerst der Antichrist kommt, oder dass die Schrecken der Trübsal beginnen. Jesus kommt wie ein Dieb in der Nacht.
        Es muss nicht sein, dass vor der Entrückung eine Verfolgung stattfinden wird. In Matthäus 24,41-42 steht, dass die Menschen zur Zeit der Entrückung ganz normal ihren Alltag leben werden. Oder in Lukas 21, 34-36, dass wir aufpassen sollen, dass wir nicht in Sünden verstrickt sind und ein berauschtes Luxusleben führen.
        Jesus sagt in \bibleverse{Luk}(21:34-36):
        \begin {bibelbox}{SCHL}{Luk}{21:34-36}
            Habt aber acht auf euch selbst, dass eure Herzen nicht beschwert werden durch Rausch und Trunkenheit und Sorgen des Lebens, und jener Tag unversehens über euch kommt. Denn wie ein Fallstrick wird er über alle kommen, die auf dem ganzen Erdboden wohnen. Darum wacht allezeit und bittet, dass ihr gewürdigt werdet, diesem allem zu entfliehen, was geschehen soll, und zu stehen vor dem Sohn des Menschen.
        \end{bibelbox}
        Es kann aber auch sein, dass für uns Bibeltreuen Christen harte Zeiten kommen. Wichtig ist, dass wir uns nicht auf diese weltlichen Dinge verlassen und vertiefen, sondern auf Jesus unseren Retter schauen.
        
        Wir haben das Wort und dort drin steht, wie wir uns auf diese Zeit vorbereiten können. Es wird entweder die Entrückung sein, oder aber unser natürlicher Tod hier auf Erde. Auch diesen Zeitpunkt weiss nur Gott. Wir wissen aber das wir zum Tod hinleben.

        Wenn es soweit ist, können wir gewiss sein, dass die, die Jesus angenommen haben, von Jesus mit offenen Armen empfangen werden. Das schönste ist, wenn wir von ihm mit den Worten: \enquote{Willkommen, du guter und treuer Knecht.} oder \enquote{Willkommen, du gute und treue Magd.} Aber das dürfen wir uns erarbeiten.
    \end{block}
    Irgendwann nach der Entrückung wird die 7-jährige Trübsal beginnen. Es wird eine schreckliche Zeit sein und das Volk Israel wird dort wieder eine Rolle spielen. Davon rede ich das nächste Mal. 
    
    Jetzt singen wir zusammen ein Lied und dann wird Thomas mit dem zweiten Teil des Tausendjährigen Reiches fortfahren. Wie wir letztes Mal gehört haben, fängt dieses Reich direkt nach der Trübsalszeit an. Dann wenn Jesus Christus in Macht und Herrlichkeit für alle sichtbar wiederkommt und die Herrschaft übernimmt.

    \lied{Alle Herren dieser Welt}.
\end{spacing}
\section{Predigt}
\textbf{Nach der Predigt}\newline
Danken für die Predigt.
\section{Abendmahl}
\begin{itemize}
    \item [] Beten für das Brot
    \item [] Beten für den Wein
\end{itemize}
\section{Abschluss}
\begin{itemize}
    \item [] \lied{Schau ich zurück}
    \item [] \beten{}
\end{itemize}
\begin{bibelbox}{SCHL}{IIKor}{13:13}
    Die Gnade des Herrn Jesus Christus und die Liebe Gottes und die Gemeinschaft des Heiligen Geistes sei mit euch allen! Amen
\end{bibelbox}
\end{document}